\chapter{Происхождение родительского действия конечномерного модуля концевого сектора}
\label{ch:v9-endpoint-finite-module-parent-origin}

\section{Дискретные данные старой конечной геометрии}

Предыдущий гейт построил минимальную архитектуру, но не её происхождение.
Старая конечная геометрия имеет фиксированные данные
\begin{equation}
 \mathfrak F_{21}=(\mathcal A_{21},H_{21},\pi_{21},J_{21},\Gamma_{21};D_{21}),
 \qquad D_{21}\in\operatorname{End}(H_{21}).
 \label{eq:v9-endpoint-parent-fixed-triple}
\end{equation}
В спектральном действии и его внутренних флуктуациях варьируется оператор
$D_{21}$, тогда как $\mathcal A_{21}$, $H_{21}$ и $\pi_{21}$ задают область
вариации \cite{v9-endpoint-parent-spectral-action,v9-endpoint-parent-dynamic-triple}:
\begin{equation}
 \delta D_{21}\in\operatorname{End}(H_{21}),
 \qquad
 \delta\dim H_{21}=0,
 \qquad
 \delta[\pi_{21}]=0.
 \label{eq:v9-endpoint-parent-variation-domain}
\end{equation}

Для трёх затронутых типов запишем кратности как
\begin{equation}
 m_{21}=(m_{0,+},m_{0,-},m_{-1,+})=(0,0,2),
 \qquad
 m_{24}=(1,1,3).
 \label{eq:v9-endpoint-parent-multiplicity-vectors}
\end{equation}
Требуемый скачок равен
\begin{equation}
 \Delta m=m_{24}-m_{21}=(1,1,1),
 \qquad \lVert\Delta m\rVert_1=3.
 \label{eq:v9-endpoint-parent-multiplicity-jump}
\end{equation}
Кратности являются целочисленными данными representation module и диаграммы
Краевского \cite{v9-endpoint-parent-krajewski}. В фиксированной компоненте
пространства полей их касательное пространство нулевое:
\begin{equation}
 T_{m_{21}}\mathbb Z_{ge0}^{3}=\{0\},
 \qquad
 \frac{\delta S_{21}}{\delta m}\ \text{не определено}.
 \label{eq:v9-endpoint-parent-zero-multiplicity-tangent}
\end{equation}

\section{Исчерпание механизмов старого родительского функционала}

Полиномиальное, коммутаторное и функциональное исчисление не выходит из
старой operator algebra:
\begin{equation}
 p(D_{21},[D_{21},\pi_{21}(a)],\Gamma_{21},J_{21})
 \in\operatorname{End}(H_{21}).
 \label{eq:v9-endpoint-parent-old-operator-closure}
\end{equation}
Внутренняя флуктуация также сохраняет carrier и его кратности:
\begin{equation}
 D_A=D_{21}+A+\epsilon'J_{21}AJ_{21}^{-1}
 \in\operatorname{End}(H_{21}),
 \qquad m(D_A)=m_{21}.
 \label{eq:v9-endpoint-parent-inner-fluctuation}
\end{equation}
Конденсат, kernel или spectral projector старого оператора может выбрать
лишь подпространство $H_{21}$ и потому не создаёт отсутствующий trivial
gauge type:
\begin{equation}
 K\le_G H_{21}
 \quad\Longrightarrow\quad
 \dim\operatorname{Hom}_G(\mathbb C_0,K)=0.
 \label{eq:v9-endpoint-parent-subspace-no-neutral}
\end{equation}

Morita/linking completion требует исходного бимодуля; Stinespring-дилатация
сохраняет различие system/environment; суммы по конечным геометриям и меры
между диаграммами в старом родительском функционале нет. Семь candidate mechanisms дают
\begin{equation}
 \begin{aligned}
 N_{\rm fixed\ parent\ module\ origin}&=0/7,\\
 \mathcal C_{\rm old}&=\{\delta D,D_A,\operatorname{Alg}(D),\ker D,\\
 &\hspace{3.5em}\text{Morita},\text{environment},\text{geometry sum}\}.
 \end{aligned}
 \label{eq:v9-endpoint-parent-candidate-ledger}
\end{equation}

\section{Условный projector-parent и его вырождение}

Чтобы увидеть минимально недостающую структуру, временно предположим, что
$H_{24}$ уже дан, и введём множество rank-three projectors
\begin{equation}
 \mathfrak P_3(H_{24})=
 \{P=P^*=P^2\in\operatorname{End}(H_{24}):\Tr P=3\}.
 \label{eq:v9-endpoint-parent-projector-space}
\end{equation}
Без типизированного ориентира естественный положительный functional
\begin{equation}
 V_0(Q)=\kappa\Tr(Q^2-Q)^2
       +\mu(\Tr Q-3)^2,
 \qquad Q=Q^*,\quad \kappa,\mu>0,
 \label{eq:v9-endpoint-parent-projector-potential}
\end{equation}
имеет множеством нулевых минимумов весь Grassmannian:
\begin{equation}
 \operatorname*{argmin}V_0=\mathfrak P_3(H_{24})
 \simeq \operatorname{Gr}_{\mathbb C}(3,24),
 \qquad \dim_{\mathbb R}=2\cdot3\cdot21=126.
 \label{eq:v9-endpoint-parent-grassmannian-degeneracy}
\end{equation}
Следовательно, rank penalty выбирает число новых линий, но не их gauge,
grading и family type.

Пусть $P_{\mathcal E}$ --- projector на требуемую
$\mathcal E_{\min}$, и введём оператор с spectral gap
\begin{equation}
 C_{\mathcal E}=I_{24}-2P_{\mathcal E},
 \qquad
 \Spec C_{\mathcal E}=\{-1^{\times3},+1^{\times21}\}.
 \label{eq:v9-endpoint-parent-target-seed}
\end{equation}
Тогда на constrained пространство конфигураций функционал
\begin{equation}
 S_C(P)=\varepsilon\Tr(C_{\mathcal E}P),
 \qquad P\in\mathfrak P_3(H_{24}),\quad\varepsilon>0
 \label{eq:v9-endpoint-parent-seeded-score}
\end{equation}
действительно имеет единственный минимум
\begin{equation}
 \operatorname*{argmin}_{P\in\mathfrak P_3(H_{24})}S_C(P)
 =\{P_{\mathcal E}\},
 \qquad \min S_C=-3\varepsilon.
 \label{eq:v9-endpoint-parent-seeded-minimum}
\end{equation}
Но этот успех тавтологичен:
\begin{equation}
 P_{\mathcal E}=\frac{I_{24}-C_{\mathcal E}}2.
 \label{eq:v9-endpoint-parent-seed-circularity}
\end{equation}
Typed seed уже содержит искомый новый module, а само определение $P$ уже
предполагает ambient carrier $H_{24}$.

Кроме того, increment-projector выделяет одну координату внутри
неприводимого family-triplet. На $H_{\rm t}\simeq\mathbb C^3$ для
$P_{\rm new}=\operatorname{diag}(1,0,0)$ и генератора, смешивающего первые
две координаты, имеем
\begin{equation}
 J_{12}=E_{12}-E_{21},
 \qquad [J_{12},P_{\rm new}]\ne0,
 \qquad \rank[J_{12},P_{\rm new}]=2.
 \label{eq:v9-endpoint-parent-family-projector-obstruction}
\end{equation}
По лемме Шура family-инвариантный projector на irreducible triplet может
быть только $0$ или $I_3$. Значит, target seed не только тавтологичен, но и
фиксирует семейную ось, если symmetry breaking не введено отдельно.

\section{Итог происхождения}

Итак, условная projector-архитектура математически работоспособна, но
старое действие не поставляет ни пространство конфигураций конечных геометрий,
ни ambient module, ни нетавтологичный selector:
\begin{equation}
 \begin{aligned}
 N_{\rm conditional\ projector\ architecture}&=5/5,\\
 N_{\rm finite\ module\ parent\ origin}&=0/5,\\
 N_{\rm raw\ physical\ slot\ closure}&=0/4.
 \end{aligned}
 \label{eq:v9-endpoint-parent-final-ledgers}
\end{equation}

\begin{theorem}[No-go fixed-parent происхождения endpoint-модуля]
Вариация спектрального действия и внутренних флуктуаций на фиксированной
конечной геометрии сохраняет размерность Hilbert carrier и целочисленные
кратности representation types. Поэтому скачок $(0,0,2)\to(1,1,3)$ не
может быть стационарным решением старого родительского функционала. Rank-three projector-parent
условно существует лишь после задания $H_{24}$; без typed seed он имеет
126-мерное многообразие минимумов, а seed с единственным минимумом прямо
кодирует $P_{\mathcal E}$ и выделяет ось внутри family-triplet.
\end{theorem}

\begin{nogocriterion}[Запрет содержащий целевую подстановку projector seed]
Оператор $C_{\mathcal E}=I-2P_{\mathcal E}$ нельзя считать объяснением
происхождения $\mathcal E_{\min}$: из него projector восстанавливается
алгебраически. Физический parent обязан сначала определить общее
пространство конфигураций конечных геометрий и нетавтологичную меру или локальный
закон перехода между его компонентами.
\end{nogocriterion}

\begin{protocol}\begin{enumerate}
\item скачок multiplicity vector: \textbf{$(1,1,1)$};
\item новых complex lines: \textbf{$3$};
\item $D$-variation меняет Hilbert module: \textbf{нет};
\item inner fluctuations меняют multiplicities: \textbf{нет};
\item fixed-parent candidate origin: \textbf{$0/7$};
\item rank-three projector architecture: \textbf{$5/5$ условно};
\item размерность вырожденного Grassmannian: \textbf{$126$};
\item target seed выбирает $P_{\mathcal E}$: \textbf{да, тавтологично};
\item increment-projector family-инвариантен: \textbf{нет};
\item physical parent-origin: \textbf{$0/5$};
\item следующий гейт: \textbf{admission общего пространство конфигураций конечных геометрий}.
\end{enumerate}\end{protocol}

Машинный сертификат сохранён в
\path{s2t_v9_endpoint_finite_module_parent_action_origin_gate_results.json}.

\begin{thebibliography}{9}
\bibitem{v9-endpoint-parent-spectral-action}
A.~H.~Chamseddine, A.~Connes,
\emph{The Spectral Action Principle}, arXiv:hep-th/9606001.
\bibitem{v9-endpoint-parent-krajewski}
T.~Krajewski, \emph{Classification of Finite Spectral Triples},
arXiv:hep-th/9701081.
\bibitem{v9-endpoint-parent-dynamic-triple}
C.~Rovelli, \emph{Spectral noncommutative geometry and quantization:
a simple example}, arXiv:gr-qc/9904029.
\end{thebibliography}